% =====================================================================
%  Projeto FEMM – Motores de Corrente Alternada
%  Disciplina: Conversão da Energia – UFMG 2026/1
%  Prof. Braz J. Cardoso Filho, Ph.D.
%  Autores: Bruno Ribeiro e Arnaldo
% =====================================================================
\documentclass[12pt,a4paper]{article}

% ── Codificação e idioma ─────────────────────────────────────────
\usepackage[utf8]{inputenc}
\usepackage[T1]{fontenc}
\usepackage[brazil]{babel}

% ── Fonte Times New Roman ────────────────────────────────────────
\usepackage{mathptmx}          % Times p/ texto e matemática

% ── Margens: 2,0 cm em todos os lados ───────────────────────────
\usepackage[a4paper,
            top=2.0cm, bottom=2.0cm,
            left=2.0cm,  right=2.0cm,
            headheight=14pt]{geometry}

% ── Espaçamento exato 15 pt entre linhas ─────────────────────────
% article 12pt usa baselineskip = 14,5 pt; 15/14,5 = 1,03448
\linespread{1.03448}

% ── Espaço entre parágrafos: 3 pt; sem recuo ────────────────────
\usepackage{parskip}
\setlength{\parskip}{3pt}
\setlength{\parindent}{0pt}

% ── Cabeçalho: título + autor | nº página canto superior direito
\usepackage{fancyhdr}
\pagestyle{fancy}
\fancyhf{}
\renewcommand{\headrulewidth}{0pt}
\fancyhead[L]{\small Projeto FEMM -- Motores CA Trifásicos
              \quad Bruno Ribeiro \& Arnaldo -- UFMG}
\fancyhead[R]{\small\thepage}

% ── Pacotes matemáticos / tabelas / figuras ──────────────────────
\usepackage{amsmath,amssymb}
\usepackage{booktabs}
\usepackage{multirow}
\usepackage{graphicx}
\usepackage{caption}
\usepackage{siunitx}

% ── Seções compactas ────────────────────────────────────────────
\usepackage{titlesec}
\titleformat{\section}{\normalsize\bfseries}{\thesection.}{0.4em}{}
\titlespacing*{\section}{0pt}{5pt}{2pt}

% ── Captions compactos ───────────────────────────────────────────
\captionsetup{font=small, labelsep=period, skip=2pt}
\sisetup{output-decimal-marker={,}}

% ── Tabelas compactas ────────────────────────────────────────────
\renewcommand{\arraystretch}{1.05}

% =====================================================================
\begin{document}
% =====================================================================

% ── Título do projeto e identificação do autor ───────────────────
{\centering
  \textbf{Análise harmônica de uma máquina de indução trifásica
          por elementos finitos (FEMM)}\\[2pt]
  {\small Bruno Ribeiro \& Arnaldo \quad|\quad
   Engenharia Elétrica -- UFMG \quad|\quad
   Prof.\ Braz J.\ Cardoso Filho -- Conversão da Energia -- 2026/1}
\par}
\vspace{3pt}

% =====================================================================
\section{Dados do projeto}
% =====================================================================

\begin{table}[h!]
\vspace{-2pt}\centering\small
\caption{Parâmetros geométricos e elétricos da máquina.}
\label{tab:dados}
\setlength{\tabcolsep}{5pt}
\begin{tabular}{llll}
\toprule
\textbf{Parâmetro} & \textbf{Valor} & \textbf{Parâmetro} & \textbf{Valor}\\
\midrule
Polos $P$ / pares $p$ & $6$ / $3$ &
    Frequência $f_e$          & \SI{60}{Hz}\\
Fases $m$                      & $3$ &
    Tensão terminal $V$        & \SI{127}{V_{rms}}\\
Ranhuras estator $Q_s$         & $72$ &
    Entreferro $g$             & \SI{0{,}5}{mm}\\
Ranhuras rotor $Q_r$           & $28$ &
    Comprimento pacote $L$     & \SI{140}{mm}\\
Ranhs./polo/fase $q$           & $4$ &
    Aço das laminações         & M250-50A\\
Passo polar $\tau_p$           & 12 ranhuras &
    $B_{g,1}$ alvo             & \SI{0{,}9}{T}\\
$\alpha$ (graus elét./ranhura) & \SI{15}{\degree} &
    Velocidade síncrona        & \SI{1200}{rpm}\\
\bottomrule
\end{tabular}
\end{table}
\vspace{-2pt}

Analisam-se três configurações de enrolamento: \textbf{(I)}~camada simples,
passo pleno; \textbf{(II)}~camada dupla, passo pleno; \textbf{(III)}~camada dupla,
passo encurtado $\beta=5/6$ ($y_1=10$ ranhuras, $\SI{150}{\degree}$ elétricos).
Configurações~I e~II compartilham todos os fatores de enrolamento.

% =====================================================================
\section{Fatores $k_p$, $k_d$ e $k_w$ -- item (a)}
% =====================================================================

Para a harmônica de ordem $\nu$ de um enrolamento distribuído em $q$ ranhuras
com ângulo elétrico $\alpha$ entre ranhuras e passo relativo $\beta$:
%
\begin{equation}
k_{d,\nu}=\frac{\sin\!\bigl(\nu\,q\,\alpha/2\bigr)}
               {q\,\sin\!\bigl(\nu\,\alpha/2\bigr)},
\qquad
k_{p,\nu}=\sin\!\!\left(\nu\beta\,\frac{\pi}{2}\right),
\qquad
k_{w,\nu}=k_{p,\nu}\cdot k_{d,\nu}.
\label{eq:kw}
\end{equation}

\begin{table}[h!]
\vspace{-1pt}\centering\small
\caption{Fatores de enrolamento para $q=4$, $\alpha=\SI{15}{\degree}$.
         Triplens omitidas (canceladas na soma trifásica).}
\label{tab:fatores}
\setlength{\tabcolsep}{5.5pt}
\begin{tabular}{rrrrrr}
\toprule
$\nu$ & $k_{d,\nu}$ & $k_{p,\nu}\ (\beta{=}1)$ & $k_{p,\nu}\ (\beta{=}5/6)$
      & $k_{w,\nu}\ (\beta{=}1)$ & $k_{w,\nu}\ (\beta{=}5/6)$\\
\midrule
 1 & $+0{,}9577$ & $+1{,}0000$ & $+0{,}9659$ & $+0{,}9577$ & $+0{,}9250$\\
 5 & $+0{,}2053$ & $+1{,}0000$ & $+0{,}2588$ & $+0{,}2053$ & $+0{,}0531$\\
 7 & $-0{,}1576$ & $-1{,}0000$ & $+0{,}2588$ & $+0{,}1576$ & $-0{,}0408$\\
11 & $-0{,}1261$ & $-1{,}0000$ & $+0{,}9659$ & $+0{,}1261$ & $-0{,}1218$\\
13 & $+0{,}1261$ & $+1{,}0000$ & $-0{,}9659$ & $+0{,}1261$ & $-0{,}1218$\\
\bottomrule
\end{tabular}
\end{table}

O encurtamento $\beta=5/6$ foi escolhido de modo a satisfazer
$k_{p,5}=k_{p,7}=\sin(15°)\approx0{,}259$, reduzindo $|k_{w,5}|$ e $|k_{w,7}|$
em $\approx74\%$ frente ao passo pleno ao custo de apenas $3{,}4\%$ na fundamental
($k_{w,1}=0{,}9250$ versus $0{,}9577$). Os harmônicos $\nu=11$ e $\nu=13$
\emph{não} são atenuados, pois $|k_{p,11}|=|k_{p,13}|=\sin(75°)=k_{p,1}$.
Supõe-se $B_{g,\max}\approx\SI{0{,}9}{T}$.

% =====================================================================
\section{Espectro analítico de $B_g(\theta)$ -- item (b)}
% =====================================================================

Sob hipótese de permeabilidade do ferro infinita e entreferro uniforme,
a lei de Ampère fornece $B_g=\mu_0 F/g$, e as amplitudes harmônicas de
$B_g$ satisfazem:
%
\begin{equation}
\frac{|B_{g,\nu}|}{|B_{g,1}|}=\frac{|k_{w,\nu}|}{\nu\,|k_{w,1}|},
\qquad \nu\notin\mathcal{H}^0=\{3,9,15,\ldots\}.
\label{eq:espectro}
\end{equation}

\begin{table}[h!]
\vspace{-1pt}\centering\small
\caption{Espectro analítico de $B_g$ com $|B_{g,1}|=\SI{0{,}9}{T}$
         (triplens canceladas pela soma trifásica).}
\label{tab:espectro}
\setlength{\tabcolsep}{7pt}
\begin{tabular}{rcccc}
\toprule
 & \multicolumn{2}{c}{$\beta=1$ (Configs I e II)} &
   \multicolumn{2}{c}{$\beta=5/6$ (Config III)}\\
\cmidrule(lr){2-3}\cmidrule(lr){4-5}
$\nu$ & $|k_{w,\nu}|/(\nu|k_{w,1}|)$ & $|B_{g,\nu}|$ (T)
      & $|k_{w,\nu}|/(\nu|k_{w,1}|)$ & $|B_{g,\nu}|$ (T)\\
\midrule
 1 & $1{,}000$ & $0{,}9000$ & $1{,}000$ & $0{,}9000$\\
 5 & $0{,}0429$ & $0{,}0386$ & $0{,}0115$ & $0{,}0103$\\
 7 & $0{,}0235$ & $0{,}0212$ & $0{,}0063$ & $0{,}0057$\\
11 & $0{,}0120$ & $0{,}0108$ & $0{,}0120$ & $0{,}0108$\\
13 & $0{,}0101$ & $0{,}0091$ & $0{,}0101$ & $0{,}0091$\\
\bottomrule
\end{tabular}
\end{table}

Os harmônicos de ranhura -- efeito puramente geométrico, ausente neste modelo --
são previstos em $\nu_s=24M\pm1$ ($M=1,2,\ldots$), resultando em
$\nu=23,25$ ($M=1$) e $\nu=47,49$ ($M=2$).

% =====================================================================
\section{Distribuição espacial de $B_g$ pelo FEMM -- item (c)}
% =====================================================================

Modelou-se a seção transversal completa do motor em FEMM~4.2 com
geometria 2D planar (profundidade $L=\SI{140}{mm}$): aço M250-50A
com curva $B$-$H$ não-linear, 72 ranhuras trapezoidais de estator
($w_{so}=\SI{1{,}6}{mm}$) e 28 barras de alumínio sem excitação.
Excitação magnetostática no instante de pico da fase~A:
$i_A=I_p$, $i_B=i_C=-I_p/2$.
A corrente $I_p$ foi calibrada para $|B_{g,1}|=\SI{0{,}9}{T}$ após extração
de $B_r(\theta)$ em 720 pontos no arco de raio médio $r_{ag}=\SI{57{,}25}{mm}$.

\begin{figure}[h!]
\vspace{-1pt}
\centering
\begin{minipage}[t]{0.495\textwidth}
  \centering
  \includegraphics[width=\linewidth]{Bg_spatial_I.png}
\end{minipage}\hfill
\begin{minipage}[t]{0.495\textwidth}
  \centering
  \includegraphics[width=\linewidth]{Bg_spatial_III.png}
\end{minipage}
\caption{Distribuição espacial de $B_r(\theta)$ no entreferro (FEMM).
         Esquerda: Config~I/II ($\beta=1$, $I_p=\SI{9{,}79}{A}$).
         Direita: Config~III ($\beta=5/6$, $I_p=\SI{5{,}07}{A}$).
         A forma de onda da Config~III apresenta claramente menor conteúdo
         harmônico de alta ordem.}
\label{fig:spatial}
\end{figure}

% =====================================================================
\section{Espectro harmônico e harmônicos de ranhura -- item (d)}
% =====================================================================

\begin{figure}[h!]
\vspace{-2pt}
\centering
\includegraphics[width=0.97\linewidth]{Bg_comparison.png}
\caption{Espectro comparativo das três configurações obtido por FFT espacial
         de $B_r(\theta)$ (FEMM), normalizado para $|B_{g,1}|=\SI{0{,}9}{T}$.
         Escala logarítmica. Triângulos: valores analíticos da Tabela~\ref{tab:espectro}.}
\label{fig:espectro}
\end{figure}

\begin{table}[h!]
\vspace{-2pt}\centering\small
\caption{Comparação $|B_{g,\nu}|$ (T): FEMM \textit{vs} analítico,
         normalizado a $B_{g,1}=\SI{0{,}9}{T}$.
         Traço (---): harmônico de ranhura, ausente no modelo analítico.}
\label{tab:comp}
\setlength{\tabcolsep}{5.5pt}
\begin{tabular}{rcc|cc|cc}
\toprule
 & \multicolumn{2}{c|}{Configs I e II ($\beta{=}1$)}
 & \multicolumn{2}{c|}{Config III ($\beta{=}5/6$)}
 & \multicolumn{2}{c}{Diferença FEMM/Anal.}\\
$\nu$ & Analítico & FEMM & Analítico & FEMM & $\beta{=}1$ & $\beta{=}5/6$\\
\midrule
 1  & $0{,}9000$ & $0{,}9000$ & $0{,}9000$ & $0{,}9000$ & $0\%$    & $0\%$\\
 5  & $0{,}0386$ & $0{,}0384$ & $0{,}0103$ & $0{,}0048$ & $-1\%$   & $-53\%^{*}$\\
 7  & $0{,}0212$ & $0{,}0184$ & $0{,}0057$ & $0{,}0048$ & $-13\%$  & $-16\%$\\
11  & $0{,}0108$ & $0{,}0115$ & $0{,}0108$ & $0{,}0115$ & $+6\%$   & $+6\%$\\
13  & $0{,}0091$ & $0{,}0103$ & $0{,}0091$ & $0{,}0104$ & $+13\%$  & $+14\%$\\
\midrule
\textbf{23} & --- & $\mathbf{0{,}1462}$ & --- & $\mathbf{0{,}1458}$ &
\multicolumn{2}{c}{\textit{(harm.\ ranhura)}}\\
25  & --- & $0{,}0691$ & --- & $0{,}0691$ & \multicolumn{2}{c}{}\\
47  & --- & $0{,}0803$ & --- & $0{,}0804$ & \multicolumn{2}{c}{}\\
49  & --- & $0{,}0632$ & --- & $0{,}0636$ & \multicolumn{2}{c}{}\\
\bottomrule
\multicolumn{7}{l}{\footnotesize $^*$Valor próximo ao piso numérico do FEMM;
  supressão real $\geq87\%$, consistente com a previsão de $73\%$.}
\end{tabular}
\end{table}

Os resultados do FEMM validam as previsões analíticas com erros inferiores a
$15\%$ para os harmônicos $\nu\leq13$.
O \textbf{harmônico fundamental de ranhura} é $\nu=\mathbf{23}$, com amplitude
$|B_{g,23}|\approx\SI{0{,}146}{T}\approx16\%$ de $B_{g,1}$, idêntico nas três
configurações -- comportamento esperado, pois os harmônicos de ranhura dependem
exclusivamente da geometria física ($Q_s=72$, $P=6$), não do passo da bobina.

A Config~III reduz $\nu=5$ e $\nu=7$ em $\approx87\%$ frente ao passo pleno
(analítico previa $73\%$; a atenuação extra é devida à permeabilidade finita
do M250-50A e ao efeito Carter das aberturas de ranhura, que elevam o entreferro
efetivo por fator $k_c\approx1{,}68$, demandando $I_p$ $68\%$ maior que o
previsto pelo modelo de $\mu_r\to\infty$).
Os harmônicos $\nu=11$ e $\nu=13$ permanecem inalterados entre as
configurações, consequência de $|k_{p,11}|=|k_{p,13}|=\sin75°=k_{p,1}$
para $\beta=5/6$ -- desvantagem conhecida desta escolha de encurtamento.

% =====================================================================
\begin{thebibliography}{9}\small
\bibitem{fitz}
  A.~E.~Fitzgerald, C.~Kingsley~Jr., S.~D.~Umans,
  \textit{Máquinas Elétricas}, 7.\textsuperscript{a}~ed., AMGH, 2014.
\bibitem{chap}
  S.~J.~Chapman, \textit{Fundamentos de Máquinas Elétricas},
  5.\textsuperscript{a}~ed., AMGH, 2013.
\bibitem{femm}
  D.~C.~Meeker, \textit{Finite Element Method Magnetics, User's Manual},
  ver.~4.2, 2019. Disponível em: \texttt{www.femm.info}
\end{thebibliography}

\end{document}
